\textbf{Phần 1. Trắc nghiệm lựa chọn}
\setcounterpageref{ex}{0}

\begin{ex} %[5D12V75]
Giá trị của biểu thức $A = 0{,}25 \times 610 \times 4 + 0{,}5 \times 390 \times 2$ là
\choice
{$900$}
{\True $1\,000$}
{$1\,100$}
{$1\,200$}
\loigiai{
Ta có:
$$A = (0{,}25 \times 4) \times 610 + (0{,}5 \times 2) \times 390 = 1 \times 610 + 1 \times 390 = 610 + 390 = 1\,000$$
\textbf{Đáp án của câu hỏi là B}
}
\end{ex}

\begin{ex} %[5D12V75]
Tìm một số thập phân biết rằng nếu dời dấu phẩy của số đó sang bên phải một chữ số thì ta được số mới lớn hơn số ban đầu là $112{,}5$ đơn vị. Số thập phân đó là
\choice
{$11{,}25$}
{\True $12{,}5$}
{$15{,}2$}
{$125$}
\loigiai{
Khi dời dấu phẩy của số thập phân sang phải một chữ số thì số đó gấp lên $10$ lần.\\
Hiệu số phần bằng nhau là:
$$10 - 1 = 9\text{ (phần)}$$
Số thập phân ban đầu là:
$$112{,}5 : 9 = 12{,}5$$
\textbf{Đáp án của câu hỏi là B}
}
\end{ex}

\begin{ex} %[5D12V75]
Một cửa hàng nhập về một chiếc xe đạp với giá gốc là $1\,600\,000$ đồng. Để đạt mức lãi $15\%$ so với giá vốn, cửa hàng cần bán chiếc xe đạp đó với giá là
\choice
{$1\,800\,000\text{ đồng}$}
{\True $1\,840\,000\text{ đồng}$}
{$1\,850\,000\text{ đồng}$}
{$1\,900\,000\text{ đồng}$}
\loigiai{
Số tiền lãi mong muốn là:
$$1\,600\,000 \times 15 : 100 = 240\,000\text{ (đồng)}$$
Giá bán chiếc xe đạp là:
$$1\,600\,000 + 240\,000 = 1\,840\,000\text{ (đồng)}$$
\textbf{Đáp án của câu hỏi là B}
}
\end{ex}

\begin{ex} %[5D12V75]
Hai thành phố cách nhau $210\text{ km}$. Lúc $7$ giờ $15$ phút, một xe máy đi từ A đến B với vận tốc $36\text{ km/h}$ và một ô tô đi từ B về A với vận tốc $48\text{ km/h}$. Hai xe gặp nhau vào thời điểm nào?
\choice
{$9$ giờ $30$ phút}
{\True $9$ giờ $45$ phút}
{$10$ giờ}
{$9$ giờ $15$ phút}
\loigiai{
Tổng vận tốc của hai xe là:
$$36 + 48 = 84\text{ (km/h)}$$
Thời gian hai xe gặp nhau là:
$$210 : 84 = 2{,}5\text{ (giờ)} = 2\text{ giờ } 30\text{ phút}$$
Hai xe gặp nhau lúc:
$$7\text{ giờ } 15\text{ phút} + 2\text{ giờ } 30\text{ phút} = 9\text{ giờ } 45\text{ phút}$$
\textbf{Đáp án của câu hỏi là B}
}
\end{ex}

\begin{ex} %[5D12V75]
Một mảnh ruộng hình thang có đáy lớn $40\text{ m}$, đáy bé $25\text{ m}$, chiều cao $20\text{ m}$. Cứ mỗi $100\text{ m}^2$ thu hoạch được $70\text{ kg}$ thóc. Hỏi cả thửa ruộng thu hoạch được bao nhiêu tạ thóc?
\choice
{$45{,}5\text{ tạ}$}
{\True $4{,}55\text{ tạ}$}
{$455\text{ tạ}$}
{$0{,}455\text{ tạ}$}
\loigiai{
Diện tích mảnh ruộng hình thang là:
$$\dfrac{(40 + 25) \times 20}{2} = 650\text{ (m}^2\text{)}$$
Khối lượng thóc thu hoạch được là:
$$650 : 100 \times 70 = 455\text{ (kg)} = 4{,}55\text{ (tạ)}$$
\textbf{Đáp án của câu hỏi là B}
}
\end{ex}

\begin{ex} %[5D12V75]
Một khối gỗ hình lập phương có diện tích toàn phần bằng $216\text{ cm}^2$. Thể tích của khối gỗ hình lập phương đó là
\choice
{$144\text{ cm}^3$}
{\True $216\text{ cm}^3$}
{$196\text{ cm}^3$}
{$256\text{ cm}^3$}
\loigiai{
Diện tích một mặt hình lập phương là:
$$216 : 6 = 36\text{ (cm}^2\text{)}$$
Vì $6 \times 6 = 36$ nên cạnh của khối gỗ hình lập phương dài $6\text{ cm}$.\\
Thể tích khối gỗ là:
$$6 \times 6 \times 6 = 216\text{ (cm}^3\text{)}$$
\textbf{Đáp án của câu hỏi là B}
}
\end{ex}

\begin{ex} %[5D12V75]
Một bể nước hình hộp chữ nhật có kích thước bên trong lòng bể: dài $2\text{ m}$, rộng $1{,}5\text{ m}$, cao $1{,}2\text{ m}$. Hiện tại bể đang chứa lượng nước bằng $65\%$ dung tích của bể. Cần đổ thêm bao nhiêu lít nước nữa để đầy bể?
\choice
{$1\,150\text{ l}$}
{\True $1\,260\text{ l}$}
{$1\,350\text{ l}$}
{$1\,400\text{ l}$}
\loigiai{
Thể tích bể nước là:
$$2 \times 1{,}5 \times 1{,}2 = 3{,}6\text{ (m}^3\text{)} = 3\,600\text{ (l)}$$
Số phần trăm lượng nước cần đổ thêm là:
$$100\% - 65\% = 35\%$$
Số lít nước cần đổ thêm là:
$$3\,600 \times 35 : 100 = 1\,260\text{ (l)}$$
\textbf{Đáp án của câu hỏi là B}
}
\end{ex}

\begin{ex} %[5D12V75]
Tìm $x$ biết: $\left(x - \dfrac{1}{3}\right) : \dfrac{3}{4} = \dfrac{8}{9}$.
\choice
{$\dfrac{2}{3}$}
{\True $1$}
{$\dfrac{4}{3}$}
{$\dfrac{5}{6}$}
\loigiai{
Ta có:
$$x - \dfrac{1}{3} = \dfrac{8}{9} \times \dfrac{3}{4} = \dfrac{2}{3}$$
$$x = \dfrac{2}{3} + \dfrac{1}{3} = 1$$
\textbf{Đáp án của câu hỏi là B}
}
\end{ex}

\begin{ex} %[5D12V75]
Nếu tăng bán kính của một hình tròn lên gấp $2$ lần thì diện tích của hình tròn đó tăng lên gấp bao nhiêu lần?
\choice
{$2\text{ lần}$}
{\True $4\text{ lần}$}
{$6\text{ lần}$}
{$8\text{ lần}$}
\loigiai{
Diện tích hình tròn ban đầu: $S_1 = r \times r \times 3{,}14$.\\
Khi bán kính tăng gấp 2 lần: $R = 2 \times r$.\\
Diện tích hình tròn mới là:
$$S_2 = (2 \times r) \times (2 \times r) \times 3{,}14 = 4 \times (r \times r \times 3{,}14) = 4 \times S_1$$
Vậy diện tích tăng lên gấp $4$ lần.\\
\textbf{Đáp án của câu hỏi là B}
}
\end{ex}

\begin{ex} %[5D12V75]
Biểu đồ quạt tròn biểu diễn cơ cấu cây ăn quả trong một trang trại gồm: Xoài chiếm $35\%$, Nhãn chiếm $30\%$, Cam chiếm $25\%$, Bưởi chiếm $10\%$. Biết số cây nhãn trong vườn là $60$ cây. Tổng số cây ăn quả của cả trang trại là
\choice
{$180\text{ cây}$}
{\True $200\text{ cây}$}
{$240\text{ cây}$}
{$300\text{ cây}$}
\loigiai{
Tổng số cây ăn quả của trang trại là:
$$60 : 30 \times 100 = 200\text{ (cây)}$$
\textbf{Đáp án của câu hỏi là B}
}
\end{ex}

\begin{ex} %[5D12V75]
Bác thợ thứ nhất làm xong một công việc trong $4$ giờ, bác thợ thứ hai làm xong công việc đó trong $6$ giờ. Nếu cả hai bác cùng làm chung thì sau $1$ giờ làm được bao nhiêu phần công việc?
\choice
{$\dfrac{1}{3}\text{ công việc}$}
{\True $\dfrac{5}{12}\text{ công việc}$}
{$\dfrac{1}{2}\text{ công việc}$}
{$\dfrac{7}{12}\text{ công việc}$}
\loigiai{
Trong 1 giờ, bác thợ thứ nhất làm được $\dfrac{1}{4}$ công việc.\\
Trong 1 giờ, bác thợ thứ hai làm được $\dfrac{1}{6}$ công việc.\\
Trong 1 giờ, cả hai bác cùng làm được:
$$\dfrac{1}{4} + \dfrac{1}{6} = \dfrac{5}{12}\text{ (công việc)}$$
\textbf{Đáp án của câu hỏi là B}
}
\end{ex}

\begin{ex} %[5D12V75]
Một khu rừng phòng hộ có diện tích $15\text{ ha } 600\text{ m}^2$. Viết số đo diện tích khu rừng đó dưới dạng số thập phân có đơn vị là héc-ta:
\choice
{$15{,}6\text{ ha}$}
{\True $15{,}06\text{ ha}$}
{$15{,}006\text{ ha}$}
{$156\text{ ha}$}
\loigiai{
Ta có: $600\text{ m}^2 = \dfrac{600}{10\,000}\text{ ha} = 0{,}06\text{ ha}$.\\
Vậy: $15\text{ ha } 600\text{ m}^2 = 15 + 0{,}06 = 15{,}06\text{ ha}$.\\
\textbf{Đáp án của câu hỏi là B}
}
\end{ex}

\textbf{Phần 2. Trắc nghiệm đúng sai}
\setcounterpageref{ex}{0}

\begin{ex} %[5D12V75]
Xét tính đúng sai của các biểu thức và khẳng định toán học sau:
\choiceTF
{\True $3\text{ m}^3 45\text{ dm}^3 = 3{,}045\text{ m}^3$.}
{\True $2\text{ giờ } 15\text{ phút} = 2{,}25\text{ giờ}$.}
{Khi cạnh của một hình lập phương tăng gấp $3$ lần thì thể tích tăng gấp $9$ lần.}
{\True Tỉ số phần trăm của $25$ và $80$ là $31{,}25\%$.}
\loigiai{
\begin{itemchoice}
    \itemch Đúng. $3\text{ m}^3 45\text{ dm}^3 = 3 + 45 : 1000 = 3{,}045\text{ m}^3$.
    \itemch Đúng. $2\text{ giờ } 15\text{ phút} = 2 + 15 : 60 = 2{,}25\text{ giờ}$.
    \itemch Sai. Cạnh tăng 3 lần thì thể tích tăng $3 \times 3 \times 3 = 27\text{ lần} \neq 9\text{ lần}$.
    \itemch Đúng. $25 : 80 = 0{,}3125 = 31{,}25\%$.
\end{itemchoice}
}
\end{ex}

\begin{ex} %[5D12V75]
Một tiệm bánh nhập về $200$ chiếc bánh với giá vốn là $15\,000\text{ đồng/chiếc}$. Xét tính đúng sai của các khẳng định sau:
\choiceTF
{\True Tổng số tiền vốn nhập bánh là $3\,000\,000\text{ đồng}$.}
{\True Cửa hàng bán $150$ chiếc đầu tiên với giá $20\,000\text{ đồng/chiếc}$ thu về $3\,000\,000\text{ đồng}$.}
{Số bánh còn lại ($50$ chiếc) được bán giảm giá $12\,000\text{ đồng/chiếc}$ thu về $500\,000\text{ đồng}$.}
{\True Tổng số tiền lãi sau khi bán hết toàn bộ $200$ chiếc bánh là $600\,000\text{ đồng}$.}
\loigiai{
\begin{itemchoice}
    \itemch Đúng. Tiền vốn: $200 \times 15\,000 = 3\,000\,000\text{ đồng}$.
    \itemch Đúng. Tiền thu 150 chiếc đầu: $150 \times 20\,000 = 3\,000\,000\text{ đồng}$.
    \itemch Sai. Tiền thu 50 chiếc sau: $50 \times 12\,000 = 600\,000\text{ đồng} \neq 500\,000\text{ đồng}$.
    \itemch Đúng. Tổng tiền thu được: $3\,000\,000 + 600\,000 = 3\,600\,000\text{ đồng}$. Tiền lãi: $3\,600\,000 - 3\,000\,000 = 600\,000\text{ đồng}$.
\end{itemchoice}
}
\end{ex}

\textbf{Phần 3. Trả lời ngắn}
\setcounterpageref{ex}{0}

\begin{ex} %[5D12V75]
Tìm một số biết rằng $45\%$ của số đó có giá trị là $135$.
\shortans{300}
\loigiai{
Số cần tìm là:
$$135 : 45 \times 100 = 300$$
\textbf{Đáp số:} $300$
}
\end{ex}

\begin{ex} %[5D12V75]
Một hình tam giác có độ dài cạnh đáy là $24\text{ cm}$ và diện tích là $144\text{ cm}^2$. Tính chiều cao tương ứng của hình tam giác theo xăng-ti-mét.
\shortans{12}
\loigiai{
Chiều cao tương ứng của hình tam giác là:
$$144 \times 2 : 24 = 12\text{ (cm)}$$
\textbf{Đáp số:} $12$
}
\end{ex}

\begin{ex} %[5D12V75]
Một ô tô chạy từ tỉnh A đến tỉnh B hết thời gian $2$ giờ $15$ phút với vận tốc không đổi là $48\text{ km/h}$. Tính độ dài quãng đường AB theo ki-lô-mét.
\shortans{108}
\loigiai{
Đổi: $2\text{ giờ } 15\text{ phút} = 2{,}25\text{ giờ}$.\\
Độ dài quãng đường AB là:
$$48 \times 2{,}25 = 108\text{ (km)}$$
\textbf{Đáp số:} $108$
}
\end{ex}

\begin{ex} %[5D12V75]
Một bể nước hình hộp chữ nhật có kích thước dài $16\text{ dm}$, rộng $12\text{ dm}$ và cao $10\text{ dm}$. Tính dung tích tối đa của bể nước theo lít.
\shortans{1920}
\loigiai{
Thể tích của bể nước là:
$$16 \times 12 \times 10 = 1\,920\text{ (dm}^3\text{)} = 1\,920\text{ (l)}$$
\textbf{Đáp số:} $1920$
}
\end{ex}

\begin{ex} %[5D12V75]
Khảo sát $250$ học sinh khối 5 của một trường tiểu học thấy có $175$ bạn thích môn bơi lội. Tính tỉ số phần trăm số học sinh thích môn bơi lội. (Chỉ ghi số).
\shortans{70}
\loigiai{
Tỉ số phần trăm số học sinh thích môn bơi lội là:
$$175 : 250 = 0{,}7 = 70\%$$
\textbf{Đáp số:} $70$
}
\end{ex}

\begin{ex} %[5D12V75]
Tính tổng của dãy số tự nhiên liên tiếp từ $1$ đến $100$: $S = 1 + 2 + 3 + \dots + 99 + 100$.
\shortans{5050}
\loigiai{
Dãy số có $100$ số hạng. Tổng của dãy số là:
$$S = \dfrac{(1 + 100) \times 100}{2} = 5\,050$$
\textbf{Đáp số:} $5050$
}
\end{ex}
